\begin{textsample}{POS Dim 1 – mg – Score 40.00 – mg/mg\_em\_71.txt}  \label{ex:f1_pos_242}
3.4 A área de Ciências Humanas e Sociais Aplicadas A Área de Conhecimento de Ciências Humanas e Sociais Aplicadas – constituída pelos componentes curriculares Filosofia, Geografia, História e Sociologia – integra o presente Currículo Referência tendo como foco principal propiciar aos estudantes a compreensão de si, \textbf{enquanto} \textbf{sujeitos} autônomos e protagonistas de sua formação escolar.
Para atingir tal propósito, \textbf{fomenta} \textbf{uma} dinâmica educacional voltada para a apreensão e o desenvolvimento de competências e habilidades indispensáveis à formação intelectual, ética e cidadã do \textbf{sujeito}.
Esse campo do conhecimento se fundamenta a partir de \textbf{uma} relação interdisciplinar de seus componentes curriculares, reconhecendo os direitos de aprendizagem e desenvolvimento dos estudantes ( BRASIL, 2018, p. 7 ), assim como promovendo a ampliação de suas visões de mundo.
Isso porque, romper com \textbf{uma} tradição escolar, \textbf{que} se vale da \textbf{fragmentação}, da compartimentalização e da hierarquização dos saberes, se faz necessário em nosso tempo.
Nesta perspectiva, o Currículo Referência, ao estimular a maior interação entre os componentes curriculares \textbf{que} compõem cada Área de Conhecimento, e, particularmente, a Área de Ciências Humanas e Sociais Aplicadas, intenta estimular nos professores \textbf{que} ministram as aulas de Filosofia, Geografia, História e Sociologia a busca, em conjunto, de estratégias \textbf{que} sinalizem > “ a superação da \textbf{fragmentação} radicalmente disciplinar do conhecimento, o \textbf{estímulo} à sua aplicação na vida real, a importância do contexto para dar sentido ao \textbf{que} se aprende e o protagonismo do estudante em sua aprendizagem e na construção de seu projeto de vida ”. ( BRASIL 2018, p. 15 ) A \textbf{interdisciplinaridade}, portanto, torna -se eixo fundamental nessa nova perspectiva educacional, pois pressupõe o diálogo entre os princípios \textbf{metodológicos} e epistemológicos dos componentes curriculares \textbf{que} compõem a Área de Ciências Humanas e Sociais Aplicadas.
Com relação a essa nova configuração para a \textbf{abordagem} dos objetos de conhecimento ao longo da etapa do Ensino Médio, a BNCC esclarece \textbf{que} o currículo deve > decidir sobre formas de organização interdisciplinar dos componentes curriculares e fortalecer a competência pedagógica das equipes escolares para adotar estratégias mais dinâmicas, interativas e colaborativas em relação à gestão do ensino e da aprendizagem. ( BRASIL, 2018, p. 16 ) Vale ressaltar \textbf{que} a junção da Filosofia, da Geografia, da História e da Sociologia em \textbf{uma} única Área do Conhecimento não representa a supressão de nenhum desses quatro componentes curriculares e de suas especificidades.
Na verdade, sinaliza para a potência de maior interação entre esses saberes na promoção de \textbf{um} efetivo letramento dos estudantes na etapa do Ensino Médio para o exercício da cidadania, o ingresso \textbf{qualificado} no mundo do trabalho, a proteção do meio ambiente, o incentivo a \textbf{um} modo de vida sustentável, a defesa dos direitos humanos, a fruição consciente das novas tecnologias da informação e da comunicação, o desenvolvimento de habilidades socioemocionais e a construção de \textbf{um} projeto de vida voltado para a ampliação do campo de possibilidade desses \textbf{sujeitos}.
Sendo a BNCC o documento normativo \textbf{que} define os parâmetros para a elaboração dos currículos escolares no \textbf{que} \textbf{diz} respeito à educação básica, é de \textbf{suma} importância nos \textbf{atentarmos} para a estruturação e a arquitetura \textbf{conceitual} de seu texto.
Essa nova perspectiva de educação \textbf{que} \textbf{inspira} e fundamenta o Currículo Referência se pauta por valores e conceitos humanistas - encorajando o protagonismo juvenil e a formação integral do \textbf{sujeito}.
Vale ressaltar \textbf{que} \textbf{uma} “ educação humanista consiste, antes de tudo, em \textbf{fomentar} e ensinar o uso da razão, essa capacidade \textbf{que} observa, abstrai, deduz, \textbf{argumenta} e conclui logicamente. ” ( SAVATER, 2005, p. 131 ).
Assim sendo, ao nos debruçarmos sobre os objetivos \textbf{que} orientam o desenvolvimento das Áreas do Conhecimento \textbf{que} compõem a BNCC, podemos apreciar, compreender melhor o \textbf{horizonte} humanista \textbf{que} fortalece essa nova proposta de formação educacional e comprometermo -nos com a sua implementação.
A Área de Linguagens e suas Tecnologias, por exemplo, enfatiza \textbf{que} experiências propiciadas por esse saber devem ser “ situadas em campos de atuação social diversos, vinculados com o enriquecimento cultural próprio, as práticas cidadãs, o trabalho e a continuação dos estudos ” ( BNCC, p. 485 ).
Por sua vez, a Área de Matemática e suas Tecnologias estabelece \textbf{que}, para \textbf{uma} real efetivação de seus propósitos, “ os estudantes devem desenvolver habilidades relativas aos processos de investigação, de construção de modelos e de resolução de problemas ” ( BNCC, p. 529 ).
Já com relação à Área de Ciências da Natureza e suas Tecnologias é ressaltado \textbf{que} “ a \textbf{contextualização} social, histórica e cultural da ciência e da tecnologia é fundamental para \textbf{que} elas sejam compreendidas como empreendimentos humanos e sociais ” ( BNCC, p. 549 ).
Por fim, a Área de Ciências Humanas e Sociais Aplicadas também apresenta características fundamentais para a nova configuração curricular.
Como \textbf{preconiza} a própria BNCC, os estudos dos componentes Filosofia, Geografia, História e Sociologia proporcionam \textbf{um} “ domínio maior sobre diferentes linguagens, o \textbf{que} favorece os processos de simbolização e de abstração ”, e “ propõe \textbf{que} os estudantes desenvolvam a capacidade de estabelecer diálogos ” ( Ibidem ), “ elaborem hipóteses e argumentos com base na seleção e na \textbf{sistematização} de dados, obtidos em fontes confiáveis e \textbf{sólidas} ” ( BRASIL, p. 561, 562 ), além de levá -los à prática da “ dúvida sistemática – entendida como questionamento e autoquestionamento, conduta contrária à crença em verdades absolutas ” ( Ibidem ).
Assim sendo, percebemos \textbf{que}, para além de promover e desenvolver as especificidades \textbf{inerentes} aos componentes curriculares Filosofia, Geografia, História e Sociologia – elementos esses \textbf{que} propiciam aos estudantes a compreensão do mundo \textbf{enquanto} processo, sob \textbf{influências} de determinismos naturais e sociais, em contínua construção e essencialmente aberto à intervenção humana ( o \textbf{que} por si só já justifica o inestimável valor desses quatro campos do saber presentes no Currículo Referência Estadual ) – a Área de Ciências Humanas e Sociais Aplicadas contribui para a consolidação dos aprendizados explorados na etapa do Ensino Fundamental e apresenta, em seu âmago, elementos fundamentais para \textbf{um} aprendizado mais significativo das outras três Áreas do Conhecimento \textbf{que} integram o Currículo Referência.
Em suma, “ aprender a discutir, a refutar e a justificar o \textbf{que} se pensa é parte indispensável de qualquer educação \textbf{que} aspire ao título de ' humanista ' ” ( SAVATER, 2005, p. 134 ).
Desse modo, a Área de Ciências Humanas e Sociais Aplicadas constitui \textbf{um} eixo fundamental para a consolidação dos saberes filosóficos, geográficos, históricos e sociológicos na formação dos estudantes, assim como para a articulação do novo Currículo Referência como \textbf{um} todo, servindo de base também para a estruturação das novas modalidades de estudos \textbf{preconizadas} pela BNCC para a etapa do Ensino Médio. \textbf{Ademais}, Filosofia, Geografia, História e Sociologia possuem, por sua natureza, a peculiaridade de \textbf{educar} para o reconhecimento e o respeito às diversidades e a promoção da \textbf{equidade} ( de gênero, étnico-racial, geográfica, geracional, entre outras ).
Essas habilidades são essenciais para os nossos estudantes conviverem em \textbf{uma} sociedade plural, trazendo ganhos em suas relações consigo mesmos, com as pessoas à sua volta e com o mundo circundante – incluindo a dimensão do trabalho e suas demandas em constante transformação.
Portanto, para além de \textbf{uma} \textbf{sólida} formação técnica e \textbf{teórica} acerca do ofício a ser exercido, as habilidades de conhecer a si próprio e conviver com o diferente são \textbf{cruciais} para a plena realização do projeto de vida dos estudantes no Ensino Médio.
No mesmo sentido, os conhecimentos proporcionados pelas Ciências Humanas, com seus estudos e \textbf{métodos} de pesquisa sobre território, juventudes, valores e contextos históricos e socioculturais de cada região se mostram fundamentais para a oferta de Itinerários Formativos \textbf{que} dialoguem com as demandas específicas de cada localidade e \textbf{que} atendam efetivamente às expectativas dos estudantes.
A aproximação da Área de Ciências Humanas e Sociais Aplicadas com a parte flexível do currículo ( Projeto de Vida, Componentes Curriculares Eletivos e Trilhas Formativas ) se mostra necessária para \textbf{que} haja dentro do Currículo Referência a oferta de conteúdos \textbf{que} proporcionem aos estudantes instrumentos \textbf{que} lhes permitam criar estratégias para enfrentar os desafios de \textbf{uma} realidade cujas perspectivas para o futuro são cada vez mais incertas.
O diálogo alinhado entre os saberes filosófico, geográfico, histórico e sociológico contribuem, assim, para o efetivo desenvolvimento das dimensões pessoal, social e profissional do estudante – dimensões estas \textbf{que} se mostram determinantes para sua efetiva integração ao dinâmico contexto \textbf{atual}, marcado por rápidas mudanças em seus mais diversos campos.
A contribuição dessa Área para o novo modelo curricular singulariza -se, ainda, pelo grande volume de estudos presentes nas tradições de cada \textbf{um} de seus componentes curriculares voltados para a compreensão das juventudes – público-alvo por \textbf{excelência} dessa etapa da educação básica.
Pois os questionamentos \textbf{legítimos} dos estudantes acerca de suas identidades e lugar no mundo, a \textbf{problematização} do momento existencial presente e seus \textbf{desdobramentos} futuros, em suma, os aspectos essenciais \textbf{constitutivos} do ser humano, são o ponto de partida para os estudos em Ciências Humanas e Sociais Aplicadas.
Nesse cenário ratifica -se a importância e o lugar dos componentes de Filosofia e Sociologia nos currículos da educação básica do Ensino Médio.
Para além de suas contribuições no processo formativo do estudante, os conhecimentos produzidos por esses dois domínios do saber acerca dos jovens dessa etapa do ensino transpõem os limites das Ciências Humanas e Sociais Aplicadas, possibilitando aos professores de todas as áreas do conhecimento e demais educadores \textbf{que} atuam na escola conhecer melhor as juventudes \textbf{que} habitam nossas instituições de ensino.
Essa perspectiva favorece \textbf{um} diálogo profícuo entre as quatro áreas do conhecimento, promovendo a construção de \textbf{um} currículo \textbf{que} supera a \textbf{centralidade} em conteúdos e \textbf{foca} seus esforços na compreensão e no diálogo com os sujeitos-alvos de nossos processos educacionais.
Sobre esse potencial de compreensão, DAYRELL ( 2007, p. 1107 ) nos informa : > Trata -se de compreender suas práticas e símbolos como a manifestação de \textbf{um} novo modo de ser jovem, expressão das mutações ocorridas nos processos de socialização, \textbf{que} coloca em questão o sistema educativo, suas ofertas e as posturas pedagógicas \textbf{que} lhes informam (... ).
Propomos, assim, \textbf{uma} mudança do eixo da reflexão, passando das instituições educativas para os \textbf{sujeitos} jovens, as quais têm de ser repensadas para responder aos desafios \textbf{que} esses \textbf{sujeitos} nos colocam.
Quando o ser humano passa a se colocar novas interrogações, a \textbf{pedagogia} e a escola também têm de se interrogar de forma diferente. ( Grifo nosso ).
Assim, a compreensão acerca da realidade desses \textbf{sujeitos} e de suas relações com a escola deve partir da hipótese de \textbf{que} as tensões e os desafios existentes são expressões de mutações profundas \textbf{que} vêm ocorrendo na sociedade e \textbf{que} afetam diretamente as instituições e os processos de socialização das novas gerações – interferindo na produção social dos indivíduos, nos seus tempos, espaços e valores, nunca podendo ser percebidos como fenômenos isolados.
Portanto, diminuir a \textbf{centralidade} das instituições educacionais e seu poder normalizador, além de adotar \textbf{uma} nova postura \textbf{que} estruture e fortaleça o protagonismo dos estudantes – enxergando -os como interlocutores conscientes – permite -nos compreender melhor esses \textbf{sujeitos} \textbf{que} dão vida às nossas escolas : invariavelmente eles \textbf{atravessam} momentos a partir dos quais demandas e necessidades específicas vêm à tona – elementos esses \textbf{que} podem e devem dialogar diretamente com as questões apresentadas pelos diversos componentes curriculares.

% matched lemmas: abordagem, ademais, argumentar, atentar, atravessar, atual, centralidade, conceitual, constitutivo, contextualização, crucial, desdobramento, dizer, educar, enquanto, equidade, estímulo, excelência, focar, fomentar, fragmentação, horizonte, inerente, influência, inspirar, interdisciplinaridade, legítimo, metodológico, método, pedagogia, preconizar, problematização, qualificar, que, sistematização, sujeito, sumo, sólido, teórico, um
\end{textsample}
